% !TITLE 声声慢·寻寻觅觅
% !AUTHOR 李清照
% !DYNASTY 两宋
% !GENRE 词
% !ORDER 12

\begin{poem}
寻寻觅觅\textsuperscript{1}，冷冷清清，凄凄惨惨戚戚\textsuperscript{2}。
乍暖还寒\textsuperscript{3}时候，最难将息\textsuperscript{4}。
三杯两盏淡酒，怎敌他\textsuperscript{5}、晚来风急。
雁过也，正伤心，却是旧时相识\textsuperscript{6}。

满地黄花堆积\textsuperscript{7}，憔悴损\textsuperscript{8}，如今有谁堪摘\textsuperscript{9}？
守着窗儿，独自怎生得黑\textsuperscript{10}？
梧桐更兼细雨，到黄昏、点点滴滴。
这次第\textsuperscript{11}，怎一个愁字了得\textsuperscript{12}！
\end{poem}

\begin{annotations}
\notetext{1}{寻寻觅觅：心神不宁，好像在找什么东西，东看看西看看。七个叠字开篇连用，在词史上是空前的创新。}
\notetext{2}{戚戚：忧伤、悲切。冷冷清清是环境，凄凄惨惨戚戚是内心——七组叠字从外到内，层层递进。}
\notetext{3}{乍暖还寒：天气刚刚转暖又变冷了。乍，忽然。这是秋天的气候特征。}
\notetext{4}{将息：调养、休息。将（jiāng），保养。天气忽冷忽热，最难保养身体。}
\notetext{5}{怎敌他：怎么抵挡得住它（指晚风）。淡酒两三杯，哪里敌得过傍晚的急风。}
\notetext{6}{旧时相识：从前认识的大雁。大雁南飞过冬——李清照是北方人，逃难到南方，看到大雁就像看到故乡的旧识。}
\notetext{7}{黄花堆积：菊花凋落堆积在地上。黄花，菊花。}
\notetext{8}{憔悴损：凋谢枯萎得不成样子。也暗指词人自己因忧愁而憔悴。}
\notetext{9}{有谁堪摘：还有谁值得采摘？堪，值得、可以。菊花败了，谁还有心思去摘呢？}
\notetext{10}{怎生得黑：怎么挨到天黑？生，语气助词。独自守在窗前，日子太难熬，天什么时候才能黑。}
\notetext{11}{次第：光景、情形。这样的境况——冷雨梧桐、黄花满地、孤身一人。}
\notetext{12}{了得：说得尽、概括得了。如此复杂的痛苦，一个“愁”字怎么说得完！}
\end{annotations}

\begin{appreciation}
《声声慢》是李清照晚年的作品，也是中国词史上叠字用得最惊人、最有名的一首之一。开篇十四个字，七组叠字连用，前无古人，后无来者。

“寻寻觅觅”——她在找什么？她自己大概也说不清。丈夫没了，家没了，国家的一半也没了，整个人像丢了魂似的。“冷冷清清”是周围的环境，“凄凄惨惨戚戚”是心里头的天气。七组叠字，从外写到内，一层比一层凉。

“乍暖还寒时候，最难将息。”深秋天气忽冷忽热，最容易生病，也最难安顿自己。“三杯两盏淡酒，怎敌他、晚来风急”——不是酒淡，是愁太重，几杯酒压不住傍晚的风。

“雁过也，正伤心，却是旧时相识。”李清照是北方人，逃难到了南方。大雁南飞，还是从前认得的那群雁——可从前是和谁一起看雁的呢。古人说雁能传书，如今信要写给谁。

下阕写白天怎么熬。“满地黄花堆积，憔悴损，如今有谁堪摘。”菊花落了一地，人也憔悴得不成样子。从前在《如梦令》里，她连一夜风雨都要为花担心；如今花真的落尽了，没有人看，也没有人摘。“守着窗儿，独自怎生得黑”——一个人守着窗子，天怎么还不黑。“梧桐更兼细雨，到黄昏、点点滴滴”，雨打在梧桐叶上，一声声的。最后一句，“这次第，怎一个愁字了得”——这样的光景，一个“愁”字哪里装得下。

对比着读很有意思：同一支笔，少女时代写“知否，知否？应是绿肥红瘦”，那是为花着急，急得可爱；晚年写“怎一个愁字了得”，是愁到说不清。中间隔着的，是国破、丧夫、流浪。所以少年人的愁未必是小事，但至少还有明天；晚年的愁，是没有明天的。

“乍暖还寒时候，最难将息”，放到现在就是一句家常话：换季的时候最容易生病。该吃饭吃饭，该睡觉睡觉，冷了添衣裳。把日子过稳当，别让这句词在你身上应验。
\end{appreciation}
